\section{Benchmark tables}\label{sec:si-benchmark}
% src: tables/coverage_tables.py (reads the stored per-cohort explorer payloads of the J and H reference runs)

Tables in this section summarise the benchmark of Section~\ref{sec:val-benchmark}. All coverage figures are for closed
95\,\% intervals, counted over the cohorts with a closed claim of that quantity (the column $n_{\mathrm{closed}}$);
Wilson intervals are 95\,\%. ``Bias'' is the median signed $\log_{10}$ error of the point estimate and ``width'' the
median closed-interval width in decades, both over the cohorts with a closed claim; ``decl.'' is the share of the
arm's cohorts for which the method made no claim for that quantity. $\kappa$ is closed on very few cohorts and is
listed for completeness.

\begin{table}[H]
\centering\small
\resizebox{\textwidth}{!}{\begin{tabular}{lrrrrrrr}
\toprule
Arm & main & bign & sharp & n-ladder & library & reference run & valid in reference run \\
\midrule
pure & 666 & 22 & 112 & 300 & 1100 & 688 & 688 \\
censor & 218 & 22 & 56 & -- & 296 & 240 & 240 \\
hazard & 218 & 22 & 56 & -- & 296 & 240 & 142 \\
window & 218 & 22 & 52 & -- & 292 & 240 & 240 \\
het & 218 & 22 & 56 & -- & 296 & 240 & 142 \\
realfull & 106 & 22 & -- & -- & 128 & 128 & 92 \\
\midrule
all & 1644 & 132 & 332 & 300 & 2408 & 1776 & 1544 \\
\bottomrule
\end{tabular}
}
\caption{Design of library v6: cohorts per arm and group. ``Reference run'' is the population of the benchmark runs
(main and big-$n$ groups); ``valid'' counts the reference-run cohorts not flagged by the validity audit.}
\label{tab:library-design}
\end{table}

\begin{table}[H]
\centering\small
\begin{tabular}{lcccc}
\toprule
 & $\eta$ & $\beta$ & $\varepsilon$ & $\kappa$ \\
\midrule
J one-sided & 0.895 & 0.710 & 0.735 & 0.534 \\
H one-sided & 0.966 & 0.980 & 0.955 & 0.969 \\
null closed & 1.000 & 1.000 & 1.000 & 1.000 \\
\bottomrule
\end{tabular}

\caption{Coverage of one-sided bounds (J and H) and of the closed interval of the data-ignoring null prior, pooled over
all arms of the reference run (not per arm). The null prior returns the whole search range, so it covers by
construction; the information is in the widths and decline rates, not in this row.}
\label{tab:onesided-null}
\end{table}

{\small
\begingroup
\captionof{table}{Closed-interval coverage, bias, width and decline rate by arm, method (J, H) and parameter on the
reference run (main and big-$n$ groups, \C{gate-v6-n} cohorts).}\label{tab:coverage-full}
\begin{longtable}{llrrrcrrr}
\toprule
Arm & Meth. & Par. & $n_{\mathrm{closed}}$ & hits & coverage [Wilson 95\%] & bias & width & decl. \\
\midrule
\endfirsthead
\toprule
Arm & Meth. & Par. & $n_{\mathrm{closed}}$ & hits & coverage [Wilson 95\%] & bias & width & decl. \\
\midrule
\endhead
\bottomrule
\endfoot
pure & J & $\eta$ & 644 & 617 & 0.958 [0.940, 0.971] & +0.05 & 0.45 & 0.000 \\
pure & H & $\eta$ & 623 & 594 & 0.953 [0.934, 0.967] & +0.05 & 0.41 & 0.000 \\
pure & J & $\beta$ & 281 & 272 & 0.968 [0.940, 0.983] & +0.27 & 3.29 & 0.004 \\
pure & H & $\beta$ & 271 & 263 & 0.970 [0.943, 0.985] & +0.28 & 3.31 & 0.054 \\
pure & J & $\varepsilon$ & 654 & 627 & 0.959 [0.941, 0.971] & +0.05 & 0.45 & 0.000 \\
pure & H & $\varepsilon$ & 637 & 609 & 0.956 [0.937, 0.969] & +0.05 & 0.45 & 0.000 \\
pure & J & $\kappa$ & 0 & 0 & -- & -- & -- & 0.919 \\
pure & H & $\kappa$ & 0 & 0 & -- & -- & -- & 0.949 \\
\midrule
censor & J & $\eta$ & 199 & 194 & 0.975 [0.943, 0.989] & +0.04 & 0.79 & 0.000 \\
censor & H & $\eta$ & 197 & 188 & 0.954 [0.915, 0.976] & +0.04 & 0.79 & 0.000 \\
censor & J & $\beta$ & 41 & 38 & 0.927 [0.806, 0.975] & +0.61 & 3.60 & 0.033 \\
censor & H & $\beta$ & 43 & 39 & 0.907 [0.784, 0.963] & +0.57 & 3.61 & 0.096 \\
censor & J & $\varepsilon$ & 201 & 196 & 0.975 [0.943, 0.989] & +0.05 & 0.96 & 0.000 \\
censor & H & $\varepsilon$ & 188 & 179 & 0.952 [0.912, 0.975] & +0.05 & 0.85 & 0.000 \\
censor & J & $\kappa$ & 0 & 0 & -- & -- & -- & 0.821 \\
censor & H & $\kappa$ & 0 & 0 & -- & -- & -- & 0.883 \\
\midrule
hazard & J & $\eta$ & 237 & 18 & 0.076 [0.049, 0.117] & +3.17 & 0.56 & 0.000 \\
hazard & H & $\eta$ & 178 & 178 & 1.000 [0.979, 1.000] & +0.03 & 0.82 & 0.000 \\
hazard & J & $\beta$ & 7 & 2 & 0.286 [0.082, 0.641] & +1.65 & 3.29 & 0.042 \\
hazard & H & $\beta$ & 20 & 20 & 1.000 [0.839, 1.000] & +0.42 & 3.80 & 0.271 \\
hazard & J & $\varepsilon$ & 209 & 11 & 0.053 [0.030, 0.092] & +3.98 & 0.38 & 0.000 \\
hazard & H & $\varepsilon$ & 171 & 170 & 0.994 [0.968, 0.999] & +0.05 & 0.99 & 0.004 \\
hazard & J & $\kappa$ & 4 & 0 & 0.000 [0.000, 0.490] & +3.63 & 1.96 & 0.171 \\
hazard & H & $\kappa$ & 0 & 0 & -- & -- & -- & 0.904 \\
\midrule
window & J & $\eta$ & 235 & 130 & 0.553 [0.489, 0.615] & +0.05 & 0.12 & 0.000 \\
window & H & $\eta$ & 181 & 181 & 1.000 [0.979, 1.000] & +0.04 & 1.11 & 0.000 \\
window & J & $\beta$ & 55 & 48 & 0.873 [0.760, 0.937] & +0.70 & 3.19 & 0.000 \\
window & H & $\beta$ & 17 & 17 & 1.000 [0.816, 1.000] & +0.61 & 4.80 & 0.229 \\
window & J & $\varepsilon$ & 234 & 77 & 0.329 [0.272, 0.392] & +0.15 & 0.17 & 0.000 \\
window & H & $\varepsilon$ & 173 & 173 & 1.000 [0.978, 1.000] & +0.07 & 1.36 & 0.004 \\
window & J & $\kappa$ & 2 & 0 & 0.000 [0.000, 0.658] & +2.82 & 1.79 & 0.800 \\
window & H & $\kappa$ & 0 & 0 & -- & -- & -- & 0.912 \\
\midrule
het & J & $\eta$ & 233 & 37 & 0.159 [0.117, 0.211] & +0.63 & 0.52 & 0.000 \\
het & H & $\eta$ & 211 & 36 & 0.171 [0.126, 0.227] & +0.47 & 0.44 & 0.004 \\
het & J & $\beta$ & 135 & 13 & 0.096 [0.057, 0.158] & +1.32 & 3.44 & 0.058 \\
het & H & $\beta$ & 166 & 24 & 0.145 [0.099, 0.206] & +1.44 & 3.53 & 0.133 \\
het & J & $\varepsilon$ & 230 & 13 & 0.057 [0.033, 0.094] & +1.09 & 0.55 & 0.000 \\
het & H & $\varepsilon$ & 203 & 18 & 0.089 [0.057, 0.136] & +0.85 & 0.50 & 0.000 \\
het & J & $\kappa$ & 0 & 0 & -- & -- & -- & 0.650 \\
het & H & $\kappa$ & 0 & 0 & -- & -- & -- & 0.867 \\
\midrule
realfull & J & $\eta$ & 123 & 54 & 0.439 [0.354, 0.527] & +1.37 & 0.41 & 0.000 \\
realfull & H & $\eta$ & 87 & 83 & 0.954 [0.888, 0.982] & +0.69 & 2.42 & 0.031 \\
realfull & J & $\beta$ & 53 & 42 & 0.792 [0.665, 0.880] & +0.68 & 3.22 & 0.016 \\
realfull & H & $\beta$ & 26 & 26 & 1.000 [0.871, 1.000] & +0.23 & 3.55 & 0.531 \\
realfull & J & $\varepsilon$ & 112 & 41 & 0.366 [0.283, 0.458] & +1.73 & 0.28 & 0.000 \\
realfull & H & $\varepsilon$ & 63 & 61 & 0.968 [0.891, 0.991] & +0.32 & 1.79 & 0.016 \\
realfull & J & $\kappa$ & 3 & 0 & 0.000 [0.000, 0.561] & +2.28 & 1.63 & 0.750 \\
realfull & H & $\kappa$ & 0 & 0 & -- & -- & -- & 0.844 \\
\end{longtable}

\endgroup
}
